% Pregunta 5. Competitividad no precio y política comercial (bloque III.3) · peso bajo
\section*{Pregunta 5. Competitividad no precio y aranceles}

\begin{enunciado}
\textbf{5.} Valore la \textbf{competitividad no precio} de la economía española. ¿Cree que la política arancelaria de la actual administración de Estados Unidos ha tenido algún impacto sobre la economía española? ¿Qué medidas de política comercial y de internacionalización recomendaría?

\smallskip
\textbf{Peso bajo} · \tiempo{13} · guía: III.3 (no precio y política comercial), B.1 (aranceles de un tercero; canales de un choque externo) · cuadro 6; gráficos 8 y 9.
\end{enunciado}

\begin{examen}
\textit{La competitividad no precio de España es sólida pero no brillante: mantiene su cuota en las exportaciones mundiales y gana peso en servicios y tecnología, aunque invierte poco en innovación. Los aranceles de Estados Unidos han reducido las ventas a ese mercado, pero su efecto directo sobre la economía es pequeño.}

\rotulo{a) La competitividad no precio}
La cuota de España en las exportaciones mundiales de bienes y servicios se mantiene en torno al 2~\% desde 2007 (cuadro~6), algo notable con el auge de los emergentes, y las exportaciones en volumen crecen un 56~\% con un tipo de cambio real casi estable frente a los socios (gráfico~8): ganan mercados por algo más que el precio. Destacan los \textbf{servicios}, cuya cuota se recupera del 1,7~\% en 2020 al 2,7~\% en 2025, y no solo por el turismo: el saldo de los servicios no turísticos pasa de cero en 2007 al 2,6~\% del PIB en 2025 (cuadro~7). Las exportaciones de alta tecnología duplican su peso (del 4,2 al 8,5~\% entre 2007 y 2022). El punto débil es el \textbf{gasto en I+D}, del 1,50~\% del PIB frente al 2,26~\% de la UE, junto con un tejido de empresas pequeñas, que limita la innovación y la internacionalización.

\rotulo{b) El impacto de los aranceles de Estados Unidos}
Las exportaciones de bienes a Estados Unidos (suma de 12 meses, gráfico~9) caen un 11~\% entre marzo de 2025 y mayo de 2026, mientras las dirigidas al mundo se mantienen: hay un \textbf{efecto directo}, pero acotado, porque Estados Unidos solo absorbe en torno al 4~\% de las exportaciones de bienes (unos 17.000 millones de euros, el 1~\% del PIB). La pérdida, unos 2.000 millones en un año, equivale a algo más de una décima del PIB. Desde junio las ventas repuntan, en parte por el efecto base frente a los meses más afectados de 2025. Más relevantes pueden ser los \textbf{efectos indirectos}: la menor demanda de los socios europeos que exportan a Estados Unidos, la desviación hacia Europa de productos que ya no entran en ese mercado y, sobre todo, la \emph{incertidumbre}, que retrasa la inversión.

\rotulo{c) Medidas}
La política comercial es competencia de la UE, de modo que las recomendaciones tienen 2 planos. En el europeo, negociar con Estados Unidos, usar la defensa comercial frente a la desviación de comercio y ampliar los acuerdos con otras regiones para \textbf{diversificar mercados}. En el nacional, \textbf{internacionalizar} más empresas (promoción, financiación y seguro de crédito a la exportación, apoyo a las pymes para que exporten de forma regular) y, en el fondo, más \textbf{I+D} y \textbf{tamaño empresarial}, que son lo que sostiene la competitividad no precio.
\end{examen}

\begin{complemento}{5.1 · Por qué se contesta así, con los cálculos}
\que{Qué se aplica.} Guía, III.3 (cuota de mercado que se mantiene sin ganar competitividad precio: competitividad no precio; recomendaciones) y B.1 (efectos de los aranceles de un tercero: directo, cadenas de valor, desviación de comercio, incertidumbre y tipo de cambio).

\begin{figuracaja}
\caption{Los canales de un choque externo (infografía de la guía)}\label{fig:canales-choque}
\grafico[0.62\linewidth]{canales-choque-externo}
\fuente{Guía del dictamen, B.1. Qué hay que ver: para los aranceles de Estados Unidos se activan sobre todo el canal comercial y el de confianza; el de materias primas es el del petróleo (pregunta~4).}
\end{figuracaja}

\que{Cálculos} (gráfico~9 y \texttt{calculos.txt}): suma de 12 meses de las exportaciones a Estados Unidos, 18.660 millones en marzo de 2025 y 16.641 en mayo de 2026 ($-10{,}8~\%$); 17.412 en agosto de 2026 ($-6{,}7~\%$ frente a marzo de 2025). Al mundo, +0,6~\% en el mismo periodo. Pérdida: unos 2.020 millones, el $0{,}12~\%$ del PIB de 2025 (1.690 miles de millones). Es una cota del efecto \emph{bruto} en valor: el valor añadido perdido es menor (las exportaciones llevan contenido importado), pero los efectos indirectos lo amplían. Peso de Estados Unidos en las exportaciones de bienes (sumas de 12 meses): 4,7~\% a principios de 2025 y 4,2~\% en mayo de 2026.

\que{Servicios no turísticos} (cuadro~7): servicios menos viajes, $2{,}9-3{,}0\approx0{,}0$ en 2007 y $6{,}8-4{,}3=2{,}6$~\% del PIB en 2025.

\que{Errores típicos.} Exagerar el impacto («los aranceles hunden las exportaciones»): el dato agregado no lo muestra. Leer la suma de 12 meses como si fuera el mes: tarda en reflejar los cambios (y el repunte interanual de julio y agosto de 2026, +20 y +39~\%, se compara con los peores meses de 2025). Recomendar aranceles de represalia nacionales: la política comercial es de la UE.

\que{Qué añadir con más tiempo.} El reparto de las exportaciones por destinos (el cuadernillo solo da Estados Unidos y el total), para medir la exposición indirecta a través de los socios europeos, y el tipo de cambio del dólar, que cambia la competitividad frente a Estados Unidos. \que{Qué recortar en el examen.} Nada: es la pregunta corta; 3 párrafos de 3 o 4 frases.
\end{complemento}
